\documentclass[11pt,a4paper]{article}
\usepackage{gtu-el-yazisi}
\begin{document}
\notbilgi{Sayısal Haberleşme}{PCM, kuantalama, hat kodları ve delta modülasyonu}{Ders-6.pdf}

\sayfa{1}
\section{Darbe kod modülasyonu (Pulse code modulation)}

\elcizimi[0.85]{sekiller/Ders-6-s1-1.png}[1]

$x_p=\mathrm{maks}\,|x(t)|$

\[ \underset{\text{kuantalayıcı adımı}}{\Delta v}=\frac{2x_p}{L}\ \longrightarrow\ \text{kuantalayıcıdaki seviye sayısı} \]

$L=2^n$ olduğunda midir

\[ N_0=N_q=\frac{\Delta v^2}{12}=\frac{x_p^2}{3L^2}\ \Big\}\ \mathrm{SNR}=\frac{S_0}{N_0}=3L^2\,\frac{P_x}{x_p^2} \]

$S_0=P_x$

Ses (Voice): 0-15 kHz \quad (300-3400 Hz önemli)

\elcizimi[0.95]{sekiller/Ders-6-s1-2.png}[1]

$R=nf_s$ bit/sn

\subsection{Düzgün olmayan kuantalama (Non uniform quantization)}

\elcizimi[0.4]{sekiller/Ders-6-s1-3.png}[1]

İşaretin genliği yani gücünün düşük olduğu yerde kuantalama aralığını \okunamadi{küçültürsek}
\[ N_0=N_q=\frac{\Delta v^2}{12}=\frac{x_p^2}{3L^2}\uparrow\ \Big\}\ \mathrm{SNR}=\frac{S_0}{N_0}=3L^2\,\frac{P_x}{x_p^2} \]

\sayfa{2}
\elcizimi[0.5]{sekiller/Ders-6-s2-1.png}[2]

küçük genlikleri büyütüyor, büyük genlikleri küçültüyor

\subsection{$\mu$ tipi sıkıştırıcı (low)}
kuzey Amerika ve Japonya

\elcizimi[0.5]{sekiller/Ders-6-s2-2.png}[2]

\[ \frac{y(t)}{y_p}=\frac{1}{\ln(1+\mu)}\,\ln\!\left(1+\mu\,\frac{|x(t)|}{x_p}\right)\mathrm{sgn}(x),\qquad 0\le\frac{|x(t)|}{x_p}\le 1 \]

SNR $=40$ dB için $\mu\ge 100$ \hfill Standart $\mu=255$

\subsection{A tipi sıkıştırıcı}
Avrupa ve diğer bölgeler

\elcizimi[0.4]{sekiller/Ders-6-s2-3.png}[2]

\[ \frac{|y(t)|}{y_p}=\begin{cases} \dfrac{A}{1+\ln A}\left(\dfrac{|x(t)|}{x_p}\right)\mathrm{sgn}(x), & 0\le\dfrac{|x(t)|}{x_p}\le\dfrac{1}{A}\\[3mm] \dfrac{1}{1+\ln A}\left(1+\ln\dfrac{A|x(t)|}{x_p}\right), & \dfrac{1}{A}\le\dfrac{|x(t)|}{x_p}\le 1 \end{cases} \]


\elcizimi[0.4]{sekiller/Ders-6-s2-4.png}[2]

Alıcı\\
Açıcı (Expander)

compander $=$ compressor $+$ expander

\sayfa{3}
$\mu$-sıkıştırıcı/açıcı kullanıldığında\\
Çıkışta: İGO (SNR)

\[ \mathrm{SNR}=\frac{S_0}{N_0}\approx\frac{3L^2}{\left[\ln(1+\mu)\right]^2},\qquad \mu^2\gg\frac{x_p^2}{P_x} \]

\elcizimi[0.55]{sekiller/Ders-6-s3-1.png}[3]

\subsection{Kodlayıcı (Encoder)}
Microelectronics Circuits, Rashid, sayfa 873-883

Digit at a time encoder

\elcizimi[0.95]{sekiller/Ders-6-s3-2.png}[3]

\[ V_o=-R_F I_F=-R_F V_{\mathrm{ref}}\left(\frac{B_3}{R_3}+\frac{B_2}{R_2}+\frac{B_1}{R_1}+\frac{B_0}{R_0}\right) \]
\[ V_o=-\frac{R_F}{R}\,V_{\mathrm{ref}}\left(2^3B_3+2^2B_2+2B_1+B_0\right) \]

\sayfa{4}
$V_{\mathrm{ref}}=5\,\mathrm{V}$ \quad $R=10R_F$ olsun

\[ V_o=-0{,}5\left(8\underset{1}{B_3}+4\underset{0}{B_2}+2\underset{1}{B_1}+\underset{1}{B_0}\right)\ \longrightarrow\ V_o=-5{,}5\ \mathrm{V} \]

\subsection{İletim hızı ve BG}

$L=2^n$ seviye $\longrightarrow$ $n\,\dfrac{\mathrm{bit}}{\text{k.örnek}}$

İletim hızı: $R_b=n\,\dfrac{\mathrm{bit}}{\mbox{\text{k.örnek}}}\cdot f_s\,\dfrac{\mbox{\text{örnek}}}{\mathrm{sn}}$

\formulkutu{R_b=n f_s\ \mathrm{bit/sn}}

\begin{teorem}
1 Hz'lik band genişliğinde, 1 saniyede 2 parça bilgi iletilir.
\begin{center}
\begin{tabular}{ccc}
\underline{\hspace{1em}Hız\hspace{1em}} & & \underline{\hspace{1em}BG\hspace{1em}}\\[2pt]
$\dfrac{2\ \text{parça}}{\text{sn}}$ & $\longrightarrow$ & 1 Hz
\end{tabular}
\end{center}
\end{teorem}

PCM için gereken min kanal BG

\[ \star\quad \mathrm{BG}=B_k=\frac{R_b}{2}=\frac{n f_s}{2}\ \mathrm{Hz} \]

\begin{ornek}
$x(t)$'nin bant genişliği $3\ \mathrm{kHz}$. $f_s$, Nyquist hızının \%33 fazlasıdır. Maks. kuantalama hatası tepe genliğinin \%0,5'idir. Kuantalanmış örnekler ikili kodlanmıştır. Gereken minimum kanal BG'ni bulunuz. Bu türden 24 işaret zaman bölmeli çoklama (ZBÇ) ile iletirse gereken min kanal BG'ni bulunuz.
\end{ornek}

\[ B=f_m=3\ \mathrm{kHz} \]
\[ f_s=f_N+\%33\,f_N=\frac{4}{3}f_N=\frac{4}{3}\times 2B=8\ \mathrm{kHz}\quad(\text{8k örnek/sn}) \]

Maks kuantalama hatası \quad $\dfrac{\Delta v}{2}=\dfrac{0{,}5}{100}\,x_p$

\[ \Delta v=\frac{2x_p}{L}\ \longrightarrow\ \frac{\Delta v}{2}=\frac{x_p}{L} \]

\[ \Longrightarrow\ \frac{\Delta v}{2}=\frac{x_p}{L}=\frac{0{,}5}{100}\,x_p\ \Longrightarrow\ \kutula{L=\frac{100}{0{,}5}=200} \]

\[ \underset{256}{2^n}\geq\underset{200}{L} \]

\[ 2^n=256=2^8\ \longrightarrow\ \kutula{n=8\ \frac{\mathrm{bit}}{\text{k.örnek}}} \]

\[ R_b=n\cdot f_s=8\times 8\mathrm{k}=64\,\mathrm{k}\,\frac{\mathrm{bit}}{\mathrm{sn}} \]
\[ B_k=\frac{R_b}{2}=32\ \mathrm{kHz} \]

\sayfa{5}
\[ R_{b,\mathrm{TDM}}=24\times R_b=1{,}536\ \mathrm{M}\,\frac{\mathrm{bit}}{\mathrm{sn}} \]
\[ B_{k,\mathrm{TDM}}=24\times B_k=0{,}768\ \mathrm{MHz} \]

\[ L=2^n\ \longrightarrow\ L^2=2^{2n} \]

\[ \frac{S_0}{N_0}=\begin{cases} 3L^2\,\dfrac{P_x}{x_p^2} & \text{sıkıştırmasız}\\[3mm] \approx 3L^2\,\dfrac{1}{\left[\ln(1+\mu)\right]^2} & \text{sıkıştırmalı},\ \mu^2\gg\dfrac{x_p^2}{P_x} \end{cases} \]

\[ \frac{S_0}{N_0}=c\cdot 2^{2n} \]

\[ c=\begin{cases} \dfrac{3P_x}{x_p^2} & \text{sıkıştırmasız (Düzgün)}\\[3mm] \dfrac{3}{\left[\ln(1+\mu)\right]^2} & \text{Düzgün olmayan} \end{cases} \]

\[ B_k=\frac{n f_s}{2}\ \Longrightarrow\ n=\frac{2B_k}{f_s}\ \Longrightarrow\ \frac{S_0}{N_0}=c\cdot 2^{2n}=c\cdot 2^{\left(4B_k/f_s\right)} \]

\[ \uparrow\quad \frac{S_0}{N_0}=c\cdot 2^{\frac{4B_k}{f_s}} \]

\vurgu[turuncu]{Bant genişliğindeki küçük bir değişim işaret gürültü oranı hızla artar.}

\begin{align*}
\left(\frac{S_0}{N_0}\right)_{\mathrm{dB}} &= 10\log_{10}\frac{S_0}{N_0} = 10\log_{10}c\,2^{2n}\\
&= \underbrace{10\log_{10}c}_{\alpha}+\underbrace{20n\log_{10}2}_{6n} = \alpha+6n\ \mathrm{dB}
\end{align*}

Bit sayısını bir arttırdığında SNR 6 dB kadar yükselir.\\
Ama bant genişliği artar

\sayfa{6}
\[ \left.\begin{array}{ll} n=8 & B_k=32\ \mathrm{kHz}\\ n=9 & B_k=36\ \mathrm{kHz}\end{array}\right\}\ 32\ \mathrm{kHz}\longrightarrow 36\ \mathrm{kHz}\quad(\%12\ \text{artış}) \]
\[ \nearrow\ \frac{9\cdot\overset{4}{8}\mathrm{k}}{2} \]

15 nisan Pazartesi\\
15:30 ara sınav (Z03, Z14, B07, B02)

\begin{ornek}
BG $=4\ \mathrm{kHz}$ olan bir $x(t)$ işareti $\mu=100$'lük sıkıştırıcı kullanan bir DKM (PCM) ile iletilecektir.\\
a) $L=64$ seviyeli kuantalayıcı kullanırsa $B_k$ ve İGO?\\
b) $L=256$ \quad '' \quad '' \quad '' \quad '' \quad ''
\end{ornek}

\[ B=4\ \mathrm{kHz}\ \longrightarrow\ f_s=f_N=2B=8\mathrm{k}\ \text{örnek/sn} \]

a) $L=64=2^6\ \longrightarrow\ n=\dfrac{6\ \mathrm{bit}}{\text{k.örnek}}\ \longrightarrow\ B_k=\dfrac{n f_s}{2}=\dfrac{6\times 8\mathrm{k}}{2}=24\ \mathrm{kHz}$

b) $\text{İGO}=\alpha+6n\ \mathrm{dB}$

\[ \alpha=10\log_{10}c=10\log_{10}\frac{3}{\left[\ln(1+\underset{100}{\mu})\right]^2}=-8{,}51\ \mathrm{dB} \]
\[ \text{İGO}=-8{,}51+6\times 6=27{,}49\ \mathrm{dB} \]

\okunamadi{b)} $L=256=2^8\ \longrightarrow\ n=8\ \mathrm{bit}/\text{k.örnek}$

$\text{İGO}=\alpha+6n=-8{,}51+6\times 8=39{,}49\ \mathrm{dB}$

$B_k=\dfrac{n f_s}{2}=\dfrac{8\cdot 8\mathrm{k}}{2}=32\ \mathrm{kHz}$

\begin{center}
\begin{tabular}{lll}
 & \underline{$L=64$ ($n=6$)} & \underline{$L=256$ ($n=8$)}\\[2pt]
$B_k=$ & 24 kHz & 32 kHz \quad (\%33 artış)\\
İGO $=$ & 27,49 dB & 39,49 dB \quad (12 dB)
\end{tabular}
\end{center}

\[ \text{İGO}(L=256)=16\cdot\text{İGO}(L=64) \]

\sayfa{7}
\elcizimi[0.9]{sekiller/Ders-6-s7-1.png}[7]

\elcizimi[0.95]{sekiller/Ders-6-s7-2.png}[7]

S 87, Fig 2.22

\subsection{Parametreler}
\begin{enumerate}[label=\arabic*)]
\item DC component
\item self clocking $\to$ manchester
\item Error detection $\to$ Duo binary
\item BW compression
\item Differential encoding
\item Noise immunity \quad $P_b(\mathrm{NRZ}) < P(\text{Unipolar RZ})$
\end{enumerate}

\sayfa{8}
\sekilnotu[8]{Basılı Figure 2.23: Çeşitli PCM dalga biçimlerinin spektral yoğunlukları (Duobinary, Delay modulation, NRZ, Dicode NRZ, Bi-phase). El yazısı eklemeler: eksen etiketi ``PSD (W/Hz)'' ve ``manchester'' eğrisinin adı.}

\[ \text{Normalized BW}\ \longrightarrow\ \mathrm{BG}=\frac{W\ (\mathrm{Hz})}{R_s=\frac{1}{T}\ \text{sembol/sn}}\quad\text{(sembol hızı)} \]

En yüksek bant genişliği manchester

her darbe bir bite karşı gelirse\\
1 sembol $=1$ bit olur

$WT<1\ \Rightarrow$ BW efficient\\
$WT>1\ \Rightarrow$ BW inefficient

\subsection{Delta modülasyonu (DM)}

\[ \Delta v=\frac{2x_p}{L}\ \longrightarrow\ N_q=\frac{\Delta v^2}{12}=\frac{3L^2}{x_p^2}\footnotemark\ \longrightarrow\ \mathrm{SNR}=3L^2\,\frac{P_x}{x_p^2} \]
\footnotetext{Claude notu: muhtemelen $N_q=x_p^2/(3L^2)$ olmalı (pay ve payda yer değiştirmiş).}

kuantalama gürültüsünü azaltmak için

\[ \underset{(n\uparrow)}{L\uparrow}\ \Rightarrow\ \Delta v\downarrow\ \Rightarrow\ N_q\downarrow\ \Rightarrow\ \text{İGO}=\frac{P_x}{N_q}\uparrow\qquad\times\qquad B_k=\frac{n f_s}{2}\uparrow \]

veya $x_p$'yi küçültebiliriz

\elcizimi[0.42]{sekiller/Ders-6-s8-1.png}[8]

Differential PCM (DPCM) $\xrightarrow{\,n=1\,}$ DM

\end{document}
